\documentclass[11pt]{article}
\usepackage[margin=1in]{geometry}
\usepackage{graphicx}
\usepackage{amsmath}
\usepackage{amssymb}
\usepackage{booktabs}
\usepackage{hyperref}
\usepackage{xcolor}
\usepackage{caption}
\usepackage{microtype}
\usepackage[numbers,sort&compress]{natbib}

% Helvetica everywhere (sans-serif body + math)
\usepackage[T1]{fontenc}
\usepackage[scaled]{helvet}
\renewcommand{\familydefault}{\sfdefault}

\hypersetup{colorlinks=true, linkcolor=blue!50!black,
            citecolor=blue!50!black, urlcolor=blue!50!black}

\title{Alignment, position-based, and hybrid models of locust collective
motion: a four-metric calibration against field trajectories}

\author{Asher Albrecht \and George Stephenson \\
Department of Computer Science, University of Colorado Boulder}

\date{CSCI-5423: Biologically Inspired Multi-Agent Systems \\ Spring 2026}

\begin{document}
\maketitle

\begin{abstract}
We test which local interaction rule reproduces the collective motion of
free-ranging locust hopper bands. We simulate the Vicsek alignment model
\citep{vicsek1995}, a position-based pull model after \citet{sayin2025},
and a sequence of hybrids that combine forward-weighted alignment,
anisotropic noise, heading inertia, and a soft preferred-spacing zone. We
compare each model to the trajectory dataset of \citet{weinburd2024}
across four metrics: polarization $\Phi$, turning-angle standard deviation,
nearest-neighbor distance, and an anisotropy index of the body-centered
neighbor density evaluated both globally and inside a near-zone radial
band. The two baseline models fail in distinct, mechanistic ways: Vicsek
matches polarization only at unrealistic noise levels and produces no
spatial structure; the simplified pull model collapses to a disordered
steady state from any initial condition, for a geometric reason that we
argue rules out an entire class of instantaneous position-based rules.
Hand-tuning a hybrid closed some gaps but introduced others. We replace
hand-tuning with the Covariance Matrix Adaptation Evolution Strategy
(CMA-ES) \citep{hansen2016}, parallelized across CPU cores, and search a
six- and seven-dimensional parameter space against a weighted
relative-error loss. The calibrated six-parameter hybrid (v4) matches
polarization within 10\%, turning-angle std and nearest-neighbor distance
within 1\%, and reproduces the field's body-centered neighbor density
both globally ($A_\text{global} \approx 0$) and inside a 1.5--3.0\,cm
near-zone band ($A_\text{near} = -0.005$ for both v4 and the field). A
seven-parameter variant adding temporal bearing integration (v5) reaches
a marginally lower scalar loss but lands in a basin where the
preferred-spacing ring carries a forward-bias artifact ($A_\text{near}
= -0.06$) that the field does not have; the seventh parameter is
unnecessary at the field's measurement scale and worsens the spatial
fit. The calibrated fit depends on periodic boundary conditions; under
reflective walls it breaks. Because real hopper bands are open systems,
we treat the calibrated parameters as a successful fit to a particular
simulation regime rather than a demonstration that this rule structure
governs field hopper-band dynamics.
\end{abstract}

\section{Introduction}

Australian plague locusts (\textit{Chortoicetes terminifera}) form
hopper bands that march coherently across kilometers of terrain, causing
agricultural damage at continental scale. Two competing theories of how
they coordinate have appeared in the literature. The first, originating
with \citet{vicsek1995} and applied to locusts by \citet{buhl2006},
proposes that each individual aligns its heading with neighbors inside a
fixed radius. Order arises from a density-driven phase transition. The
second, advocated by \citet{sayin2025} on the basis of field, laboratory,
and virtual-reality experiments with desert locusts (\textit{Schistocerca
gregaria}), proposes that individuals are attracted to the positions of
forward neighbors. The interaction is geometric rather than directional,
and \citet{sayin2025} argue that bearing information is integrated over
time by a neural ring-attractor circuit, rather than read out
instantaneously. We test models inspired by both theories against
\textit{C.~terminifera} trajectory data; the species mismatch is a caveat
worth noting, although the two species share the gregarious-phase
marching behavior the theories aim to explain.

These theories make different predictions about the spatial structure of
the swarm and about the frame-to-frame heading dynamics of individuals.
Until recently they could not be tested against individual-level field
data. \citet{weinburd2024} released the first such dataset: roughly 3.3
million position records and 20{,}000 trajectories from four hopper bands,
sampled at 25\,Hz. We use this dataset as the ground truth.

We ask which local rule, or which combination of rules, reproduces the
field measurements. Our contribution is threefold. First, we show that
both pure alignment and pure pull fail in characteristic ways. The pull
failure is sharp: starting from perfect alignment ($\Phi = 1$), the model
collapses to the disordered steady state within a single timestep, which
we trace to a geometric symmetry argument that applies to any
instantaneous, reflection-symmetric, position-based rule. Second, we
replace hand-tuning of the hybrid with derivative-free optimization:
CMA-ES searches the parameter space against a multi-metric scalar loss
and finds a six-parameter hybrid (v4) that matches the three scalar
metrics. Third, we introduce a near-zone anisotropy metric that the
optimizer did not see, and use it to validate the calibrated models
against the field's body-centered spatial structure. The near-zone
metric separates v4 from a seven-parameter variant (v5) that has lower
scalar loss but a forward-biased preferred-spacing ring; this
out-of-sample test is what justifies treating v4 as the configuration
that fits the field.

\section{Background and related work}

The Vicsek model \citep{vicsek1995} is a minimal self-propelled particle
system: at each time step, every agent updates its heading to the
arithmetic mean of neighbor headings inside an interaction radius, plus
uniform angular noise of width $\eta$. \citet{buhl2006} fit this model to
laboratory locust experiments and reported a phase transition from
disorder to coherent marching as density increased. \citet{bazazi2008}
proposed an alternative mechanism in which cannibalistic interactions
drive forward motion: locusts in a marching band tend to bite one
another, and individuals risk being bitten from behind, so contact and
visual approach from behind elevate the propensity to march. Reduction
of an individual's capacity to detect rear approach (via abdominal
denervation) reduces its marching probability. Neither study had access
to free-ranging field trajectories.

\citet{weinburd2024} computed body-centered neighbor density maps from
their hopper-band data and observed an asymmetric distribution
concentrated in the near zone (within roughly 3\,cm of the focal), with
more density behind the focal than ahead. They argued for directional
interactions but stopped short of a specific mechanistic model.

\citet{sayin2025} fit a position-based model to field, laboratory, and
virtual-reality experiments. In their pull formulation, each agent steers
toward a weighted centroid of forward neighbors, with weights proportional
to $\cos^2(\beta/2)$ for relative bearing $\beta$. They argued that
bearing information is integrated over time by a neural ring-attractor
circuit and that this temporal integration is necessary for the mechanism
to produce coherent motion. Our v5 hybrid (Section~\ref{sec:methods})
incorporates this temporal-integration component to test whether it is
required to fit the Weinburd data.

The Couzin model \citep{couzin2002} combines repulsion, alignment, and
attraction in nested zones. We borrow the soft preferred-spacing idea
from this literature for our hybrid extensions.

CMA-ES \citep{hansen2016} is a derivative-free evolutionary optimizer
that adapts a multivariate Gaussian sampling distribution as it explores.
It is well suited to noisy black-box objectives in low-dimensional
continuous spaces. We use the \texttt{pycma} implementation
\citep{pycma}.

\section{Methods}
\label{sec:methods}

\subsection{Field metrics}

We compute four metrics from the Weinburd dataset and apply the same
functions to every simulated trajectory:

\begin{enumerate}
\item Global polarization $\Phi(t) = \left| \frac{1}{N(t)} \sum_i
  e^{i\theta_i(t)} \right|$, averaged over time after burn-in.
\item Turning-angle standard deviation: per-step heading change
  $\Delta\theta_i(t) = \theta_i(t+1) - \theta_i(t)$, wrapped to
  $[-\pi, \pi]$, pooled across agents and time.
\item Median nearest-neighbor distance (NND).
\item Body-centered neighbor density. For each focal individual we
  rotate the relative position of every neighbor inside a 10\,cm radius
  into the focal's body frame (heading along $+y$) and accumulate a
  $50 \times 50$ histogram. Figure~\ref{fig:fig1} shows the field
  distribution.
\end{enumerate}

From the density map we derive scalar anisotropy indices
$A = (\rho_\text{rear} - \rho_\text{front}) / (\rho_\text{rear} +
\rho_\text{front})$, where $\rho_\text{rear}$ and $\rho_\text{front}$ are
mean densities in the lower and upper halves of the body-centered
histogram. By this convention $A > 0$ indicates more density behind the
focal than ahead, consistent with the qualitative direction reported by
\citet{weinburd2024}. We compute $A$ at two scales: $A_\text{global}$
averaged over the full half-disk out to 10\,cm, and $A_\text{near}$
restricted to a 1.5--3.0\,cm radial band. The two-scale split matters
because the anisotropy \citet{weinburd2024} report is concentrated in the
near zone, so a near-isotropic value of $A_\text{global}$ is compatible
with non-trivial near-zone structure that the global metric does not
resolve. Section~\ref{sec:spatial} uses both indices: $A_\text{global}$
as a coarse first check, $A_\text{near}$ as the more discriminating
spatial test. Crucially, the CMA-ES loss only sees the three scalar
metrics; both anisotropy indices are computed post hoc, which makes
$A_\text{near}$ an out-of-sample validation of the calibrated models.

\begin{figure}[t]
\centering
\includegraphics[width=0.55\linewidth]{figures/fig1_field_density.png}
\caption{Body-centered neighbor density from the \citet{weinburd2024}
hopper-band dataset, walking and hopping individuals combined (six clips,
$6.86 \times 10^6$ neighbor counts). The focal heads up. The distribution
is close to isotropic at the global level
($A_\text{global} = -0.0008$), with a self-exclusion hole at the origin
and a near-zone anisotropy concentrated within $\sim 3$\,cm of the focal
($A_\text{near} = -0.005$). We use this map and its scalar summaries as
field-fitting targets alongside the three scalar metrics.}
\label{fig:fig1}
\end{figure}

\subsection{Models}

The Vicsek model implements standard alignment dynamics with periodic
boundaries. Each agent sets its next heading to
$\theta_i(t{+}1) = \arg\!\sum_{j \in \mathcal{N}_i}
e^{i\theta_j} + \eta \xi_i$, where $\mathcal{N}_i$ is the set of
neighbors within $r_\text{int} = 7\,$cm and $\xi_i \sim
\mathcal{U}(-1/2, 1/2)$. Speed is constant and a hard repulsion override
applies at $r_\text{rep} = 1.5\,$cm.

The pull model has each agent steer toward the centroid of forward
neighbors weighted by $w(\beta) = \cos^2(\beta/2)$, following
\citet{sayin2025}, but \emph{without} the temporal-integration component.
The noise term and the repulsion override match the Vicsek
implementation. We test this simplified per-frame version explicitly to
isolate the effect of temporal integration; v5 below restores it.

Hybrid v3 (anisotropic Vicsek) extends the alignment rule along two
axes. Forward neighbors (relative bearing $|\beta| < \pi/3$) get weight
$1 + \alpha_\text{aniso}$ and rear neighbors $1 - \alpha_\text{aniso}$.
Effective noise scales with the fraction of neighbors ahead:
$\eta_\text{eff}(i) = \eta_\text{base} (1 - \lambda f_\text{forward}(i))$.
Speed varies with local order between $v_\text{min} = 2.7\,$cm/s and
$v_\text{max} = 11.76\,$cm/s.

Hybrid v4 adds two mechanisms motivated by the v3 failure analysis
(Section~\ref{sec:results}). The heading-inertia term has each agent
commit a fraction $\mu \in (0, 1]$ of the way to its new desired
heading $\theta^*$ via exponential smoothing in unit-vector space:
$\theta_{i}(t{+}1) = \arg\!\left[(1{-}\mu) e^{i\theta_i(t)} + \mu
e^{i\theta^*}\right]$. This narrows the per-step turning distribution.
The preferred-spacing zone applies in the band $r_\text{rep} < d <
r_\text{pref}$: the agent mixes a radial-push direction into its
desired heading with weight $k_\text{pref}$.

Hybrid v5 adds temporal bearing integration after \citet{sayin2025}'s
ring-attractor analogy. Each agent maintains an exponential moving
average of the alignment-derived target unit vector, $T_\text{ema}
\leftarrow (1-\rho) T_\text{ema} + \rho T_\text{instantaneous}$. The
desired heading derives from $T_\text{ema}$ rather than the
instantaneous target. Setting $\rho = 1$ disables the integrator and
recovers v4.

\subsection{Calibration}

We calibrate the hybrid models with CMA-ES against a weighted
relative-error loss
\begin{align}
L = w_\Phi \frac{|\Phi_\text{sim} - \Phi_\text{field}|}{\Phi_\text{field}}
  + w_\sigma \frac{|\sigma_\text{sim} - \sigma_\text{field}|}{\sigma_\text{field}}
  + w_\text{NND} \frac{|\text{NND}_\text{sim} - \text{NND}_\text{field}|}{\text{NND}_\text{field}},
\end{align}
with weights $w_\Phi = 1$, $w_\sigma = w_\text{NND} = 2$. The two
failure-mode metrics from the baseline experiments
(Section~\ref{sec:results}) get double weight; polarization is already
close in v3. This weighting is an experimenter choice. We test its
sensitivity in Section~\ref{sec:discussion} by re-running v4 with equal
weights. Each evaluation runs five seeds of a 2000-step simulation with
500-step burn-in and returns the ensemble mean.

The optimization parameter space is six-dimensional for v4 and
seven-dimensional for v5: $\eta_\text{base} \in [0.1, 1.5]$,
$\lambda \in [0, 1]$, $\alpha_\text{aniso} \in [0, 1]$,
$\mu \in [0.1, 1.0]$, $k_\text{pref} \in [0, 0.7]$,
$r_\text{pref} \in [1.6, 5.0]$\,cm, and (v5 only)
$\rho_\text{temporal} \in [0.05, 1.0]$. Following the \texttt{pycma}
guidance \citep{pycma}, we map every parameter linearly to the internal
range $[0, 10]$ so a single $\sigma_0$ has uniform meaning across
dimensions.

We run CMA-ES with population size 20, $\sigma_0 = 2.0$, and 40--50
generations. The 20 candidate evaluations within each generation are
distributed across CPU cores using \texttt{multiprocessing.Pool}; BLAS
threading inside workers is pinned to one to avoid oversubscription. On
a 32-thread Intel i9-13900K the full v4 optimization completes in
about 13 minutes (gen$_0$ in 19 seconds, end-to-end speedup of about
$11{\times}$ over the serial implementation).

\section{Results}
\label{sec:results}

\subsection{The two baselines fail in distinct ways}

After calibrating $\eta$ to match field polarization, the Vicsek model
hits $\Phi = 0.824$ (target 0.820) but produces a turning-angle
distribution more than twice as wide as the field
(0.595\,rad vs 0.276\,rad, 116\% relative error) and undercounts NND by
33\%. The model achieves order, but at the cost of pathological
heading noise that compensates for the absence of any heading-stabilizing
mechanism beyond raw alignment.

The pull model fails differently, and more sharply. Across five
structural variants (varying cone half-angle from $45^\circ$ to
$180^\circ$, pull strength from 0.1 to 150, and noise from 0.1 to
2.9\,rad), polarization remained at $\Phi \approx 0.06$ in every case,
with a turning-angle distribution near $2\,$rad (essentially saturated).
The failure is not a tuning failure: it persists across the entire
parameter sweep we tested.

The failure is also not a bootstrapping failure. Initialized from
perfect order ($\Phi_0 = 1.000$, headings drawn with std $0.01\,$rad
across five seeds), polarization drops 96\% within a single timestep
and converges to the same disordered steady state as random
initialization (mean $0.0372$ vs $0.0373$, both with seed std $\lesssim
10^{-3}$). The model destroys existing order; it does not merely fail to
generate it.

The cause is geometric. In a disordered swarm, neighbor bearings are
approximately uniformly distributed around the focal individual. The
forward-weighting function $\cos^2(\beta/2)$ is symmetric under the
reflection $\beta \to -\beta$. Any rule whose weighting is symmetric
under that reflection produces, in expectation over an isotropic bearing
distribution, zero net lateral steering signal. The forward component
is non-zero, so each agent maintains its current heading; but there is
no mechanism to bias the population toward a common direction. Small
stochastic deviations from isotropy provide a finite-$N$ signal in
principle, but heading noise randomizes each agent's direction faster
than these deviations can accumulate into coherent motion. The result is
a steady state at $\Phi \approx 0.037$ that closely matches the
finite-$N$ baseline expected for $N=150$ agents moving independently
with no effective coupling.

This argument rules out an entire class of instantaneous,
reflection-symmetric, position-based steering rules as candidate
mechanisms for sustaining order in a disordered swarm. It does not rule
out the full \citet{sayin2025} proposal, which adds temporal integration
via a ring-attractor circuit. The argument is in fact consistent with
their emphasis on the ring attractor: temporal integration is what
their formulation requires precisely because the per-frame steering rule
on its own cannot do the work. Hybrid v5 (Section~\ref{sec:cma-closes})
tests whether adding temporal integration to a hybrid alignment-plus-pull
model improves the fit to the Weinburd data.

Hybrid v3 (anisotropic Vicsek with hand-tuned $\eta = 1.1$, $\lambda =
0.8$, $\alpha = 0.6$) hits polarization within 0.5\% but still misses
turning-angle std by 88\% and NND by 45\%. Forward weighting tightens
the alignment, which lowers $\eta$ enough to bring polarization in at a
sensible noise level, but introduces no mechanism that limits per-step
heading change or sets a preferred spacing.

\subsection{CMA-ES closes the scalar gaps}
\label{sec:cma-closes}

\begin{figure}[t]
\centering
\includegraphics[width=0.85\linewidth]{figures/fig2_cma_convergence.png}
\caption{CMA-ES convergence curves. Both v4 (six-dimensional) and v5
(seven-dimensional) reduce the weighted loss by an order of magnitude
over 40--50 generations. The v5 run plateaus around generation 25 before
a late breakout finds a slightly better basin. As Section~\ref{sec:spatial}
shows, the basin v5 escaped into is worse on the out-of-sample spatial
metric.}
\label{fig:fig2}
\end{figure}

Figure~\ref{fig:fig2} shows the loss trajectories. The v4 run reaches a
final loss of 0.099. The calibrated ensemble means are
$\Phi = 0.898$ (9.6\% above the field target), $\sigma = 0.276$\,rad,
and NND median $= 3.90$\,cm. The two failure modes that motivated
heading inertia and preferred spacing close to within seed-to-seed
variation: the seed std on $\sigma$ is approximately $0.005$\,rad
(roughly 2\% of the field value) and on NND is approximately $0.04$\,cm
(roughly 1\% of the field value). We deliberately do not report
percentage errors smaller than seed noise. Polarization remains
slightly elevated; the 9.6\% residual is a direct consequence of our
weighting choice (Section~\ref{sec:discussion}).

The v5 run reaches a final loss of 0.085, but the converged
$\rho_\text{temporal} = 0.95$ is near the no-effect end of its allowed
range (recall $\rho = 1$ disables the integrator). The lower scalar
loss does not come from the integration mechanism. It comes from the
optimizer finding a different basin: $\alpha_\text{aniso}$ saturates at
$1.0$ and $\lambda$ drops to $0.17$, while the other parameters move
only slightly. The two CMA-ES-calibrated runs separate further once the
spatial structure is examined directly (Section~\ref{sec:spatial}): v4
reproduces the field's near-zone density, v5 does not.

Figure~\ref{fig:fig3} summarizes per-metric performance for every model
we ran. Numerical values are listed in Table~\ref{tab:scorecard}.

\begin{figure}[t]
\centering
\includegraphics[width=\linewidth]{figures/fig3_scorecard.png}
\caption{Field-fit scorecard across all models on the three scalar
metrics. Dashed lines indicate field targets. Vicsek and Pull miss
multiple metrics. Hybrid v3 closes polarization but not the others.
CMA-ES-calibrated v4 and v5 hit all three scalar targets within seed
noise (under periodic boundary conditions, see
Section~\ref{sec:boundary}). The body-centered spatial structure
separates v4 from v5; see Figure~\ref{fig:fig4}.}
\label{fig:fig3}
\end{figure}

\begin{table}[t]
\centering
\small
\caption{Field-fit scorecard. Relative error in parentheses. Field targets:
$\Phi = 0.820$, $\sigma = 0.276$\,rad, NND $= 3.893$\,cm. $\pm$ values on
$\Phi$ are the across-seed standard deviation of the per-seed mean (10
seeds for the baseline runs, 5 seeds for the CMA-ES-calibrated runs).
Errors on $\sigma$ and NND for the calibrated hybrids fall below the
seed-to-seed standard deviation; we report them as $<$1\% for
completeness but do not interpret sub-seed-noise residuals.}
\label{tab:scorecard}
\begin{tabular}{lrrr}
\toprule
Model       & $\Phi$ \;$\pm\,\sigma_\text{seed}$       & $\sigma$ (rad)     & NND (cm)              \\
\midrule
Vicsek      & $0.824 \pm 0.019$ \;(0.5\%)              & 0.595 (116\%)      & 2.61 (33\%)           \\
Pull        & $0.058 \pm 0.001$ \;(93\%)               & 1.999 (624\%)      & 1.32 (66\%)           \\
Hybrid v3   & $0.890 \pm 0.025$ \;(8.5\%)              & 0.520 (88\%)       & 2.15 (45\%)           \\
Hybrid v4   & $0.898 \pm 0.027$ \;(9.6\%)              & 0.276 ($<$1\%)     & 3.90 ($<$1\%)         \\
Hybrid v5   & $0.889 \pm 0.012$ \;(8.4\%)              & 0.276 ($<$1\%)     & 3.89 ($<$1\%)         \\
\bottomrule
\end{tabular}
\end{table}

\subsection{Spatial structure: v4 matches, v5 distorts}
\label{sec:spatial}

The body-centered neighbor density map is the test that motivated the
hybrid extensions in the first place. \citet{weinburd2024} reported
anisotropic neighbor distributions in their data, and our hybrid v3
appeared to reproduce a visible asymmetry in the inner ring of its
density map.

The global anisotropy index is essentially zero in the field
($A_\text{global} = -0.0008$). v3, v4, and v5 all reproduce the global
value to within $|A| < 0.001$ (Figure~\ref{fig:fig4}, panel titles).
Match on $A_\text{global}$ alone is therefore a weak test: all the
calibrated models pass, but they pass by all being near-isotropic at the
global scale rather than by reproducing structured signal.

The near-zone metric $A_\text{near}$, restricted to a 1.5--3.0\,cm
radial band, separates the calibrated hybrids. The field value is
$-0.005$, marginally favoring rear-side density. v3 returns $-0.003$
and v4 returns $-0.005$, both close to the field. v5's value is
$-0.061$, an order of magnitude larger, with substantially more density
ahead of the focal than behind. The preferred-spacing ring at v5's
optimum sits closer to the focal in front than at the rear, at radii
where the field shows no asymmetry. The white dotted circles in
Figure~\ref{fig:fig4} mark the band.

\begin{figure}[t]
\centering
\includegraphics[width=\linewidth]{figures/fig4_spatial_check.png}
\caption{Body-centered neighbor density for the field and the three
hybrid models, on a shared color scale. The global anisotropy index is
within $|A| < 0.001$ of zero in all four. The near-zone index, computed
only inside the white dotted band ([1.5, 3.0]\,cm), separates the v5
ring artifact ($A_\text{near} = -0.061$) from the near-flat field
($-0.005$), v3 ($-0.003$), and v4 ($-0.005$). v4 fits both scalar and
near-zone metrics; v5 trades the spatial fit for a marginally lower
scalar loss.}
\label{fig:fig4}
\end{figure}

The CMA-ES loss did not include $A_\text{near}$, so the optimizer could
place near-zone structure anywhere that did not break the three scalar
metrics. v4 found a basin in which the preferred-spacing ring sits
symmetrically around the focal ($\alpha_\text{aniso} = 0.80$,
$\lambda = 0.52$). v5, with the extra $\rho_\text{temporal}$ dimension,
found a different basin: $\alpha_\text{aniso}$ saturates at $0.99$ and
$\lambda$ drops to $0.17$, which sharpens the ring and biases it
forward. v5's scalar loss is 0.085 versus v4's 0.099, so the optimizer
prefers v5; on the out-of-sample spatial metric the trade goes the
other way. The seventh parameter is unnecessary at the field's
measurement scale, as Section~\ref{sec:cma-closes} already noted, and
in our calibration it actively worsens the spatial fit. v4 is the
configuration that matches all four metrics; v5 is a worse fit dressed
up as a slightly better one.

The ``forward void'' claim attached to earlier v3 figures was an
overinterpretation of a near-zone feature that is small in the field
($A_\text{near} = -0.005$). v3 reproduces the field's near-zone structure
for the same reason it fails the scalar metrics: forward-weighted
alignment alone is enough for spatial structure but not for tight
heading control or NND.

\subsection{Boundary-condition dependence}
\label{sec:boundary}

We tested the calibrated v4 model under reflective walls instead of
periodic boundaries, holding all other parameters fixed. The fit broke:
turning-angle std jumped to 0.493\,rad (78\% off) and NND dropped to
3.47\,cm (11\% off). Polarization moved closer to target, from 0.898 to
0.752 (8.3\% off). The periodic image structure is doing real physical
work in the calibration. Reflective walls perturb headings on contact,
inflating the per-step turning distribution, and they compress densities
away from the boundary, lowering NND.

This is a substantive limitation, not a footnote. Real hopper bands are
open systems with front and rear asymmetry, not periodic boxes. The
calibrated parameters in Table~\ref{tab:scorecard} are conditional on a
boundary regime that does not match the system we are modeling. The same
hybrid mechanism, calibrated under different boundary conditions, would
converge on different parameters and almost certainly different fits.
We have demonstrated that with six free parameters, CMA-ES can match the
Weinburd scalar metrics and the near-zone anisotropy under periodic
boundaries; we have \emph{not} demonstrated that this rule structure
governs the dynamics of free-ranging locust hopper bands. Establishing
that would require calibration under realistic open-boundary conditions,
which we did not have time to implement. We discuss possible approaches
in Section~\ref{sec:discussion}.

\section{Discussion}
\label{sec:discussion}

The pull-model failure result and its geometric explanation stand
independently of the calibration story. They establish that no
instantaneous, reflection-symmetric, position-based steering rule can
sustain order in a disordered swarm, regardless of cone angle, pull
strength, or noise level, and they point directly at temporal integration
(or some other symmetry-breaking mechanism) as the missing ingredient.
This is consistent with \citet{sayin2025}'s mechanistic emphasis on the
ring-attractor circuit, and it tightens the claim: the ring attractor
is not just one possible implementation of pull dynamics but a
computationally necessary component of any pull-class mechanism that
generates order.

We spent several weeks hand-tuning a hybrid model and producing
scorecards that captured one metric well at a time. CMA-ES, with parallel
evaluation across cores, runs end-to-end in about 13 minutes and finds
parameter combinations that hand-tuning missed. For noisy black-box
objectives in this dimensionality, derivative-free evolutionary search
is the right default. We should have used it sooner.

Adding temporal bearing integration ($\rho_\text{temporal}$) gave the
optimizer a seventh degree of freedom that it drove near its no-effect
limit ($\rho = 0.95$). On the scalar loss alone the change looks benign
and even slightly helpful, but the extra dimension shifted the
converged basin: $\alpha_\text{aniso}$ saturated at 1.0 in v5 against
v4's 0.80, and that re-shaped the preferred-spacing ring into the
forward-biased configuration that breaks the near-zone fit
(Section~\ref{sec:spatial}). Its geometrically established necessity
(Section~\ref{sec:results}) for any \emph{pure} pull-class model
remains. What our v5 result shows is narrower: at the granularity of
the \citet{weinburd2024} dataset, and within an alignment-dominated
hybrid that already breaks reflection symmetry through forward-weighted
coupling, the additional temporal-integration term is unnecessary, and
in our calibration harmful. A measurement that distinguishes
instantaneous from time-integrated bearing estimation would be needed
to test the ring-attractor mechanism directly; neither our scalar nor
our spatial metrics do.

The remaining gap is polarization. The v4 and v5 calibrations sit at
$\Phi \approx 0.89$ versus a field target of $0.82$. We re-ran v4 with
equal weights ($w_\Phi = w_\sigma = w_\text{NND} = 1$) to test whether
the gap was an artifact of our prioritized weighting. The equal-weight
optimum lands at $\Phi = 0.883$, $\sigma = 0.277$\,rad,
NND $= 3.92$\,cm, with $\lambda$ dropping from 0.52 to 0.24 and
$\alpha_\text{aniso}$ saturating at 0.99. The polarization gap narrows
by 0.015 and the other two metrics do not break. v4's polarization
floor is therefore close to 0.88 regardless of weighting, which we
read as a structural property of the model at this density rather than
a tunable parameter.

The v3 anisotropy reframing in Section~\ref{sec:spatial} is the most
significant correction relative to our earlier work on this problem.
The visible v3 asymmetry that motivated the hybrid extensions was real
but near-zone; our initial $A_\text{global}$ index did not capture it;
and matching on $A_\text{global}$ alone is matching at zero rather than
matching on a structured signal. Introducing $A_\text{near}$ as an
out-of-sample test changes the picture: v4 passes it, v5 fails it, and
the calibrated v4 becomes a model that matches all four metrics rather
than three.

We have not run an open-boundary calibration. The reflective-wall
stress test indicates the trade-offs differ. Three plausible follow-ups
are: a flux-boundary scheme that despawns leaving agents and respawns at
the rear; calibration under reflective walls with the loss reweighted;
or a sliding-window simulation centered on the swarm centroid. Of
these, we think the sliding-window approach is most defensible because
it reproduces the actual measurement geometry of the
\citet{weinburd2024} field recordings, but we did not have time to
implement any.

\section{Project contributions}

\textbf{Asher Albrecht} built the field-data ingestion pipeline from the
\citet{weinburd2024} \texttt{.mat} files, computed the four baseline
metrics on the field data (Section 3.1), implemented the Vicsek and Pull
baseline models, ran the pull-model order-destruction experiment and
geometric analysis (Section~\ref{sec:results}), and led the hybrid v3
anisotropic-Vicsek design and the $\eta$/$\lambda$ parameter sweep
figures.

\textbf{George Stephenson} extended the hybrid model with heading inertia,
preferred spacing, and temporal bearing integration (v4 and v5 in
\texttt{hybrid\_v4.py} and \texttt{hybrid\_v5.py}), built the parallel
CMA-ES calibration framework (\texttt{calibrate\_cma.py},
\texttt{calibrate\_cma\_v5.py}, \texttt{validate\_calibrated.py}),
implemented the global and near-zone spatial-anisotropy validation
against the field density map, and produced the post-calibration
figures.

Both authors contributed to the experimental design, the slide deck,
and this paper.

\section*{Code and data}

The simulation code, calibration scripts, CMA-ES log files, and figure
generators are at \url{https://github.com/}\textit{(repo to be added on
submission)}. The Weinburd dataset is available from the original
authors at \citet{weinburd2024}.

\bibliographystyle{plainnat}
\bibliography{refs}

\end{document}
